\chapter{Related Work}

\section{Distributed Reinforcement Learning}

Early work on distributed RL established foundational architectures for scaling training. Nair et al.~\cite{nair2015} introduced massively parallel methods for deep RL using parallel actors generating behavior, parallel learners trained from stored experience, and distributed neural networks, achieving order-of-magnitude speedups on Atari games. Building on this, Espeholt et al.~\cite{impala2018} proposed IMPALA, which decouples acting from learning using importance sampling for off-policy corrections, scaling to thousands of machines while achieving state-of-the-art results on multi-task benchmarks. Horgan et al.~\cite{apex2018} introduced Ape-X, separating data collection from learning through distributed prioritized experience replay, where multiple actors generate experience stored in a shared replay buffer accessed by learners.

\section{Communication-Efficient Distributed Training}

Reducing communication overhead in distributed training has been extensively studied. Chen et al.~\cite{chen2018} developed communication-efficient policy gradient methods specifically for distributed RL, proposing gradient compression techniques that reduce communication costs while maintaining convergence guarantees. In the supervised learning domain, Lin et al.~\cite{dgc2018} discovered that 99.9\% of gradient exchange in distributed SGD is redundant, proposing Deep Gradient Compression (DGC) that achieves 270--600x bandwidth reduction through momentum correction and local gradient clipping.

\section{Pipeline Parallelism for Large Models}

Pipeline parallelism has emerged as a critical technique for training models exceeding single-device memory. Huang et al.~\cite{gpipe2019} presented GPipe, which partitions neural networks across multiple accelerators and pipelines mini-batches through stages, introducing micro-batching and gradient accumulation to achieve near-linear speedup while training models 25x larger than single-accelerator capacity. Narayanan et al.~\cite{pipedream2019} proposed PipeDream, which automatically partitions DNN models and schedules computation to minimize pipeline bubbles, maintaining multiple parameter versions to enable weight updates without stalling. For transformer models, Shoeybi et al.~\cite{megatron2019} developed Megatron-LM with efficient intra-layer model parallelism, achieving 76\% scaling efficiency on 512 GPUs for 8.3 billion parameter models. Rajbhandari et al.~\cite{zero2020} introduced ZeRO, eliminating redundant storage of optimizer states, gradients, and parameters across data-parallel processes, enabling training of 170 billion parameter models.

\section{Resource-Constrained Environments}

Recent work addresses distributed learning in resource-constrained settings. Lim et al.~\cite{lim2020} applied federated learning principles to RL for IoT devices, enabling multiple agents to collaboratively learn optimal control policies while keeping data local, addressing device-specific dynamics variations. Feng et al.~\cite{feng2023} designed IoTSL, an efficient split learning system for resource-constrained IoT devices, optimizing the split learning paradigm to handle non-IID data distribution and reduce communication requirements specific to IoT constraints.

Our work differs from prior approaches by specifically addressing the communication bottleneck in pipeline-parallel RL training through selective gradient accumulation. While existing gradient compression techniques focus on compressing all gradients uniformly or through quantization, our approach accumulates gradients over time and transmits only high-magnitude values above configurable percentile thresholds, tailored to the unique characteristics of RL training with repeated epochs.
